\section{Deneysel kurulum}\label{sec:deney-kurulum}

\subsection{Aşama 1 ve Aşama 2}\label{sec:deney-onkayit}

Aşama~2'nin bütün kararları ilk koşudan önce \path{PROTOCOL.json} dosyasında sabitlendi ve dosyanın
SHA-256 özeti kaydedildi (\texttt{cce83992\ldots}). % E791
Birincil set, projede daha önce kullanılmamış 38 VLSI problemidir (VLSI-38; $n=1.376$--$3.954$,
10 tohum, 10~sn). % E797
Dört birincil hipotez GNN'yi rastgele ve kNN hedeflemeyle iki geometride karşılaştırır: Wilcoxon
testleri tek Holm ailesinde düzeltilir, ``fark yok'' sonucu $\pm0{,}10$ puanlık sınırla iki tek
yönlü testle (TOST) sınanır. % E792
Tablo~\ref{tab:v2-sets} ikincil setleri verir. Bellek baskısıyla iki kez kesilen koşu, tamamlanmış
kayıtlara dokunulmadan sürdürüldü.

% OTOMATIK URETIM -- evidence/scripts/build_v2.py
\begin{table}[!htbp]
\centering\small\setlength{\tabcolsep}{3pt}
\caption{Hızlı ILS kampanyasının deney setleri. Her (problem, tohum) için altı strateji aynı başlangıç turundan ve aynı kick rastgele sayı tohumundan başlar.}
\label{tab:v2-sets}
\begin{tabularx}{\textwidth}{@{}>{\raggedright\arraybackslash}Xcccccl@{}}
\toprule
Set & Problem & $n$ & Tohum & Bütçe [sn] & Referans & Rol \\
\midrule
VLSI-38 (birincil) & 38 & 1376--3954 & 10 & 10 & Optimum/BKS & Ön-kayıtlı birincil \\
VLSI-20 (önceki) & 20 & 131--3293 & 10 & 10 & Optimum/BKS & İkincil \\
TSPLIB-72 & 72 & 51--2392 & 5 & 10 & Optimum/BKS & İkincil \\
Sentetik-100 (1 sn) & 100 & 81--160 & 5 & 1 & LKH-3 (30 sn) & İkincil \\
Sentetik-1000 & 40 & 1000--1000 & 5 & 10 & LKH-3 (30 sn) & İkincil \\
\bottomrule
\end{tabularx}
\end{table}


Birincil sonuçtan sonra eklenen kâhin hedefleme, sürekli tazelenen GNN, 60~sn bütçe ve iki MLP
deneyleri keşifseldir.

\subsection{Aşama 3: geliştirme ve seçim}\label{sec:deney-asama3}

Aşama~3'ün yöntem ve ayar seçimleri yalnız geliştirme setlerinde yapıldı: VLSI-38, Aşama~1'in 20
VLSI problemi (VLSI-20) ve Sentetik-1000 (5 tohum, 10~sn). Bu koşularda kNN ve fazla-uzunluk
VLSI-38'de rastgeleden iyiydi ve birbirinden ayrışmadı. % E738, E739, E740
Çevrimiçi model VLSI-38'de eşleşik sabit kolundan ($+0{,}210$ puan), diğer iki sette fazla-uzunluk
vekilinden kötüydü ($+0{,}157$ ve $+0{,}148$; ikincisi soğuk başlatılmış sürüm). % E742, E745, E747
Ardından kNN, fazla-uzunluk ve GNN için $f\in\{0{,}02;0{,}04;0{,}06;0{,}10\}$ ve
$p\in\{0{,}5;0{,}75;0{,}9\}$ ızgarası VLSI-38 ile Sentetik-1000'de 3 tohumla koşuldu. Seçim kuralı
ızgara sonuçlarından önce sapma kaydına yazıldı: iki setin ortalama gap'ini en küçükleyen hücre
seçilir; varsayılan kural en iyiye 0,01 puandan yakınsa korunur. Kural kNN için 0,04/0,75,
fazla-uzunluk için 0,06/0,5, GNN için 0,04/0,5'i seçti (Şekil~\ref{fig:v3-pool-grid}). % E750, E751, E752
Kazanç-MLP dahil diğer sinyaller ayarlanmadı.

\subsection{Aşama 3: dondurma, zaman damgası ve doğrulama seti}\label{sec:deney-dondurma}

Protokol \path{PROTOCOL_v2.json} dosyasında donduruldu; dosya hipotezleri, analizi ve kod, model ve
veri dosyalarının SHA-256 özetlerini içerir, kendi özeti \texttt{96eb1d9e\ldots8085}'tir. % E584
Kod, doğrulama setini protokol donmadan çalıştırmayı reddeder. Protokolün öngördüğü Open Science
Framework yazar için erişilemez olduğundan bu sapma kayda yazıldı; özet 12:51'de
OpenTimestamps\footnote{\url{https://opentimestamps.org}} takvim sunucularına gönderildi ve holdout2
koşusu hemen ardından başladı. Bitcoin tasdiki (blok 976258), koşu 540 görevin 384'ündeyken ve
sonuçlara bakılmadan alındı. % E585, E964
Bu blok zamanı, protokolün sonuçlar analiz edilmeden önce var olduğunu ve sonradan
değiştirilmediğini gösteren, üçüncü kişilerce bağımsız doğrulanabilen ve kurcalamaya dayanıklı üst
sınırdır; koşudan önceki sıralama yerel sapma kaydına dayanır.

Holdout2, 54 problemden oluşur: 6 Waterloo VLSI, 3 National ve 3 TSPLIB problemi ile 14 uniform,
14 kümeli ve 14 karışık sentetik problem ($n=1.689$--$9.882$). % E586, E591--E598
TSPLIB problemleri Aşama~1'de kullanılmış, ancak hiçbir Aşama~2 ya da Aşama~3 koşusunda yer
almamıştır. Sekiz strateji 10 tohum ve 10~sn ile koşuldu (toplam 4.320 koşu). % E587, E588, E590
Bir (problem, tohum) çiftinin stratejileri aynı işçide ardışık çalışır ve vuruş tohumunu da
paylaşır. Dört işçi Intel Core i5-9300H, Windows~11 x64, Python~3.13 ve Numba~0.68
ile çalıştı; bütçe duvar-saati olduğundan vuruş sayıları donanıma bağlıdır
(Tablo~\ref{tab:v3-mechanism}). Hiçbir koşu dışlanmadı.

\subsection{Aşama 3: hipotezler ve istatistik}\label{sec:deney-istatistik}

İstatistik birimi problemdir: gap önce 10 tohum üzerinden ortalanır ($n=54$). Dört doğrulayıcı
hipotez sabit sırayla, her biri ancak öncekiler reddedildiyse $\alpha=0{,}05$ düzeyinde sınanır.
V1 çevrimiçi skorun sabit doğrusal ön-bilgiden kötü, V2 vekilin rastgeleden iyi olduğunu öne sürer
(iki taraflı Wilcoxon \citep{wilcoxon_1945}). V3 ve V4, vekilin GNN'ye ve kazanç-MLP'ye
$\delta=0{,}05$ puandan fazla kötü olmadığını $d-\delta$ üzerinde tek yönlü Wilcoxon testiyle sınar
\citep{schuirmann_1987_tost, lakens_2017_equivalence}. % E589
Sınır, geliştirmede gözlenen en küçük anlamlı hedefleme etkisinin (fazla-uzunluk$-$rastgele,
$-0{,}070$ puan) altındadır; daha dar bir sınır, problem başı farkların standart sapması 0,131'e
ulaştığı için güçsüz kalırdı. % E739, E740
Öte yandan sınır, GNN'nin rastgeleye karşı Aşama~2'deki etkisiyle ($-0{,}052$ puan) aynı büyüklüktedir;
bu yüzden V3 tek başına GNN kazancının korunduğunu göstermez, V2 ve keşifsel üstünlük
karşılaştırmasıyla birlikte okunmalıdır. % E803
İkincil aile beş keşifsel üstünlük karşılaştırmasıdır (kNN$-$rastgele, fazla-uzunluk$-$kNN,
GNN$-$kNN, kazanç-MLP$-$üyelik-MLP, fazla-uzunluk$-$GNN) ve Holm ile düzeltilir \citep{holm_1979}.
Kayıtlı analiz betiği `:' içeren strateji adlarını ayrıştıramadığından aile, aynı istatistik
işlevlerini çağıran ayrı bir betikle (\path{build_v3.py}) hesaplandı. Güven aralıkları 10.000
tekrarlı bootstrap'la bulunur.
